\section{第~\ref{sec:sleep}~章　睡眠}

睡眠的模式、重放和突触的缩小已通过神经元记录、微透析和电子显微镜测量；睡眠的时间安排和慢摆动由本书的模型描述，而睡眠的目的主要还是假说。

{\footnotesize
\setlength{\tabcolsep}{3pt}\renewcommand{\arraystretch}{1.15}
\begin{longtable}{>{\raggedright}p{0.5cm}>{\raggedright}p{4.0cm}>{\raggedright}p{5.6cm}>{\raggedright}p{2.3cm}>{\raggedright\arraybackslash}p{1.7cm}}
\toprule
编号 & 论断 & 实验及测量内容 & 来源 & 状态 \\
\midrule\endfirsthead
\toprule
编号 & 论断 & 实验及测量内容 & 来源 & 状态 \\
\midrule\endhead
1 & 睡眠中皮层神经元并不沉默；改变的是工作模式 & 大鼠皮层神经元的长时间记录：慢波睡眠中频率不为零，工作期与静默期交替；长时间清醒后，各状态下的频率都更高，睡眠后则更低 & \cite{Vyazovskiy2009} & 已测量 \\
2 & \nt{ACET} 在白天和快速眼动睡眠中高，在慢波睡眠中低（表~\ref{tab:sleepmodes}） & 自由活动的猫，皮层微透析：清醒时乙酰胆碱的释放比慢波睡眠高 $100$--$175\%$，快速眼动睡眠中约与安静清醒时持平 & \cite{Marrosu1995} & 已测量 \\
3 & \nt{NORA} 和 \nt{SERO} 在慢波睡眠中安静下来，在快速眼动睡眠中几乎沉默 & 按睡眠阶段记录大鼠的单个 \nt{NORA} 神经元和猫的 \nt{SERO} 神经元：频率在清醒时最高，慢波睡眠中较低，快速眼动睡眠中几乎为零 & \cite{AstonJonesBloom1981,McGintyHarper1976} & 已测量 \\
4 & 快速眼动睡眠中肌肉经 \nt{GABA} 和 \nt{GLYC} 被抑制关闭 & 在大鼠中测量咀嚼肌 \nt{ACET} 神经元的活动，并把阻断剂直接施加到它们上：只有两类 \nt{GABA} 受体和 \nt{GLYC} 受体都被阻断时，麻痹才解除 & \cite{BrooksPeever2012} & 已测量 \\
5 & 睡眠压力 $S$ 是带两个阈值的电容器~\eqref{eq:sleeppressure}，$\tau_r=18.2$~h，$\tau_d=4.2$~h & 双过程模型；时间常数由脑电图慢波功率随清醒增长、随睡眠下降的方式得出，阈值的形状由自发醒来的时间得出 & \cite{Borbely1982,Daan1984} & 理论 \\
6 & 机器自己进入 $16.1$~h 清醒和 $8.0$~h 睡眠（图~\ref{fig:twoprocess}） & 按~\eqref{eq:sleeppressure}~加上摆动阈值计算，脚本 \texttt{models/sleep\_models.py}：清醒 $16.1$~h，睡眠 $7.9$--$8.0$~h & --- & 模型 \\
7 & $40$~h 不睡后 $S=0.90$，但睡眠只持续 $8.8$~h：欠债靠深度而不是长度偿还 & 模型：\texttt{models/sleep\_models.py} 给出 $S=0.90$ 和 $8.8$~h。实验：八个人 $40.5$~h 不睡后，第一个恢复夜中脑电图慢波功率显著高于平时 & \cite{Borbely1981} & 模型 \\
8 & 物理上 $S$ 就是腺苷 & 猫的微透析：在一个控制清醒的区域中，细胞外腺苷在清醒期间上升，睡眠剥夺时继续上升，睡眠中下降。$S$ 仅仅是腺苷这一点尚未证明 & \cite{PorkkaHeiskanen1997,Dunwiddie2001} & 假说 \\
9 & 慢波睡眠中皮层摆动：工作几百 ms——静默，每秒不到一次（$<1$~Hz，麻醉下常为 $0.3$--$0.4$~Hz） & 麻醉猫的细胞内记录：$0.8$--$1.5$~s 的去极化与长时间的超极化交替，节律 $<1$~Hz，最常见为 $0.3$--$0.4$~Hz；自然睡眠中也能看到同样的工作与静默交替（第~1 行） & \cite{Steriade1993,Vyazovskiy2009} & 已测量 \\
10 & 摆动来自带反馈和疲劳的群体~\eqref{eq:updown}；$b=5$ 时周期 $1.1$~s（模型值），约 $35\%$ 的时间工作；$b=1$ 时平稳工作（图~\ref{fig:updown}） & 计算，脚本 \texttt{models/sleep\_models.py}：周期 $1.11$~s，工作时间比例 $0.35$（按整数个周期平均）；$b=1$ 时工作比例为 $1.0$。模型的周期短于实测值（第~9 行） & --- & 模型 \\
11 & \nt{ACET} 关闭摆动，使皮层转入平稳工作 & 刺激猫的 \nt{ACET} 神经元核团，阻断了 $79\%$ 皮层神经元的慢节律，使其转入平稳工作，主要是长时间超极化消失所致。乙酰胆碱抑制产生疲劳的慢钾电流 & \cite{SteriadeAmzica1993,McCormick1992} & 已测量 \\
12 & $7$--$15$~Hz 的节律爆发由皮层和中继核团之间的 \nt{GLUT}--\nt{GABA} 环路产生 & 在雪貂视觉中继核团的脑片中，爆发由中继细胞对来自相邻 \nt{GABA} 核团的抑制所作的反跳簇放电产生，而该核团又由中继细胞自己兴奋 & \cite{vonKrosigk1993,SteriadeMcCormickSejnowski1993} & 已测量 \\
13 & 白天的重放是快进的：海马中约快 $20$ 倍，皮层中快 $6$--$7$ 倍 & 慢波睡眠中，海马神经元的重复序列在时间上被压缩。跑道奔跑后，位置神经元的序列以同样的顺序在 $\sim100$~ms 的短爆发中重复，压缩约 $20$ 倍；皮层中压缩 $6$--$7$ 倍 & \cite{Nadasdy1999,LeeWilson2002,Euston2007} & 已测量 \\
14 & 增强重放与皮层的协调可改善记忆（因果关系） & 大鼠在训练后接受与重放锁定的刺激，以加强重放与皮层慢波和节律爆发的联系：第二天记忆良好，对照组处于随机水平 & \cite{Maingret2016} & 已测量 \\
15 & 重放的目的是对慢速记忆进行交错训练 & 互补学习系统理论和人工网络上的计算：顺序学习会破坏旧知识，交错学习则不会。没有直接的脑实验表明重放恰恰为此所需 & \cite{McClelland1995} & 假说 \\
16 & 睡眠后突触按大小成比例缩小，大的几乎不变 & 对小鼠皮层 $6920$ 个突触的三维电子显微镜观察：睡眠后接触面积比清醒后小 $\sim18\%$，缩小与大小成比例，大而稳定的突触不变 & \cite{deVivo2017} & 已测量 \\
17 & 睡眠的主要任务是权重的重新归一化 & 突触稳态假说；第~1 和第~16 行与之一致，但不能证明这是主要功能 & \cite{TononiCirelli2014} & 假说 \\
18 & 快速眼动睡眠是对世界模型的台架测试 & 模式（表~\ref{tab:sleepmodes}）已测量；网络在运行世界模型是理论推测，没有实验 & \cite{Hobson2009} & 假说 \\
19 & 睡眠中大脑的冲洗：问题悬而未决 & 一项实验：睡眠中细胞间隙宽 $60\%$，清除更快。另一项采用不同测量方法的实验：睡眠中和麻醉下清除明显更慢。结果相互矛盾 & \cite{Xie2013,Miao2024} & 假说 \\
20 & 局部睡眠：清醒动物的皮层区域会短暂陷入静默 & 大鼠长时间清醒后，个别皮层区域的神经元短暂沉默，就像在慢波睡眠中一样，局部脑电图中出现慢波；此时大鼠在任务中更容易出错 & \cite{Vyazovskiy2011} & 已测量 \\
\bottomrule
\end{longtable}}
